\section{Design theory and calculations}

\subsection{Membership in a polygon}

The engine decides whether a point is inside a polygon with the ray casting method and the even-odd rule: a
ray is cast from the point and the number of polygon sides it crosses is counted; if the number is odd, the
point is inside. The vertices are given in order and the last one is joined to the first. In a polygon that
crosses itself, the region covered twice counts as outside.

A point that lies exactly on the boundary is classified with a half-open rule: in an axis-aligned rectangle,
the south side, the west side and the south-west corner count as inside, and the rest as outside. The result
on the exact boundary must not be used as data.

\subsection{Distance to the boundary}

The distance shown is the distance to the nearest side, not to the nearest vertex. To compute it, the engine
projects the coordinates onto a local plane centered on the point:

\begin{equation}
x = (\lambda - \lambda_0)\,\frac{\pi}{180}\,R\cos\varphi_0, \qquad
y = (\varphi - \varphi_0)\,\frac{\pi}{180}\,R,
\end{equation}

with $R = 6371$~km, so that one degree of latitude is about 111\,195~m. The distance is the minimum, over
all sides, of the distance from the origin to the segment. It is an unsigned distance: membership in the
polygon is decided separately.

The error of this flat projection with respect to the haversine distance is documented in the engine as
below 0.2~\% up to 85° of latitude, and it grows toward the poles. It is not used near the poles, nor with
polygons that cross the antimeridian (longitude $\pm180^\circ$), because longitudes are subtracted as plain
numbers.

\subsection{A single call that does not confuse errors with ``outside''}

The membership functions return a boolean, so they cannot tell ``outside'' from ``invalid input''. That is why
the application uses \texttt{geofence\_check()}, which in a single call returns whether the point is inside
and the distance to the boundary, and returns an error code if the fix is not valid (coordinates out of range
or not numbers, an empty polygon or an engine that was not started). An invalid fix is never interpreted as
leaving the fence.

\subsection{Square fence around a point}

For a center $(\varphi_0, \lambda_0)$ and a half-side $h$, the corners are computed with

\begin{equation}
\Delta\varphi = \frac{h}{111\,195}, \qquad \Delta\lambda = \frac{\Delta\varphi}{\cos\varphi_0},
\end{equation}

and the corners are $(\varphi_0 \mp \Delta\varphi,\ \lambda_0 \mp \Delta\lambda)$, in the order south-west,
south-east, north-east and north-west. Longitude is adjusted with the cosine of the latitude because a degree
of longitude is shorter away from the equator. The construction is rejected if the latitude of the center
exceeds 85° or if the square would cross the antimeridian.

\subsection{From distance to color}

With a safe distance $S = 20$~m and a distance $d$ to the boundary, with the point inside, the color is

\begin{equation}
t = \min\!\left(\frac{d}{S},\,1\right), \qquad
R = \operatorname{round}\bigl(255\,(1-t)\bigr), \qquad
G = \operatorname{round}\bigl(255\,t\bigr), \qquad B = 0.
\end{equation}

At $d \ge S$ the color is pure green, and at $d = 0$, steady red. Outside the fence, or on any invalid input
(a non-finite or negative distance, or $S \le 0$), the result is blinking red, the fail-safe state.

\subsection{Time budget}

The cost of \texttt{geofence\_check()} is linear in the number $N$ of vertices, because it walks every side once
to decide membership and once more to compute the distance. Measured on the board (see the results section),
the time per call fits

\begin{equation}
T(N) \approx \dBenchFitOffsetUs~\mu\text{s} + \dBenchFitSlopeUs~\mu\text{s}\cdot N,
\end{equation}

so the 50~ms limit of SR-12 would be reached near \dBenchLimitVertices{} vertices. The engine uses
double-precision numbers; since the floating-point unit of the Cortex-M33 is single precision, those
operations probably run in software, and that is probably where the cost per vertex comes from. This is an
inference from the compiler flags, not a measurement.

\subsection{Resolution of the measuring clock}

The kernel cycle counter runs at \dClockHz~Hz, so each reading has a resolution of \dTickUs~$\mu$s. A call can
last less than that, so a single call is not timed: a batch of 40 is timed and the total is divided by 40.
